% !TeX program = xelatex
\documentclass[10.5pt]{exam}
\usepackage{hyperref}
\usepackage[a4paper, total={6in, 10in}]{geometry}
\usepackage{caption}
\usepackage{float}
\usepackage{amsmath}
\usepackage{amssymb}
\usepackage{mathtools}
\usepackage{listings}
\usepackage{graphicx}
\usepackage{setspace}
\usepackage[table]{xcolor}
\usepackage{booktabs}
\usepackage{multirow}
\usepackage{array}
\usepackage{xepersian}
\settextfont{XB Niloofar.ttf}

\definecolor{tableblue}{RGB}{68,114,196}
\definecolor{lightblue}{RGB}{217,226,243}

\hypersetup{
    colorlinks=true,
    linkcolor=blue!60!black,
    urlcolor=blue!70!black
}

\lstset{
    basicstyle=\ttfamily\small,
    breaklines=true,
    frame=single,
    numbers=left,
    numberstyle=\tiny,
    keywordstyle=\color{blue},
    commentstyle=\color{gray},
    stringstyle=\color{teal},
    showstringspaces=false,
    columns=fullflexible
}

\begin{document}
\pagestyle{head}
\firstpageheader{}{}{}
\runningheader{صفحه \thepage\ از \numpages}{}{}
\runningheadrule

\begin{tabular}{p{.75\textwidth} l}
    \multirow{2}{*}{\includegraphics[scale=0.2]{UILOGO.png}} &                           \\ \\
                                                             & \textbf{علوم کامپیوتر}    \\
                                                             & \textbf{نیم‌سال دوم ۰۴-۰۵} \\
                                                             & \textbf{هوش مصنوعی}       \\ \\
    \textbf{دانشکده ریاضی و آمار}                            &                           \\
\end{tabular}\\
\noindent\rule{\textwidth}{0.5pt}

\vspace{0.1cm}

\textbf{زهرا چراغی، محمد ملائی، فاطمه سعیدی}

\noindent\rule{\textwidth}{0.5pt}

\section*{گزارش پروژه منطق مرتبه اول: حل بازی مین‌روب}

\textbf{لینک گیت‌هاب پروژه:} \lr{\url{https://github.com/Elham-ch/AI-Projects}}

\section{معرفی مسئله}
در این پروژه، بازی
\lr{Minesweeper}
به صورت یک مسئله استنتاج منطقی مدل شده است. صفحه بازی از خانه‌هایی تشکیل می‌شود که ممکن است مین داشته باشند یا امن باشند. با آشکار کردن یک خانه امن، عددی بین صفر تا هشت دیده می‌شود که تعداد مین‌های مجاور آن خانه را نشان می‌دهد. عامل باید فقط با استفاده از اطلاعاتی که تا همان لحظه آشکار شده‌اند، خانه‌های امن و مین‌ها را تشخیص دهد.

محیط بازی در فایل
\lr{\texttt{src/minesweeper.py}}
قرار دارد و کلاس
\lr{\texttt{MineSweeper}}
وظیفه تولید نقشه، نگهداری وضعیت خانه‌ها، اجرای عمل‌های
\lr{\texttt{reveal}}
و
\lr{\texttt{flag}}
و نمایش گرافیکی بازی را بر عهده دارد. نقطه شروع، خانه مرکزی صفحه است. سازنده محیط یک ناحیه سه در سه را در اطراف مرکز از مین خالی نگه می‌دارد؛ بنابراین مقدار خانه شروع همیشه صفر است.

در حل خودکار، گزینه
\lr{\texttt{auto\_flood\_fill=False}}
استفاده شده است. در نتیجه حتی اگر یک خانه صفر آشکار شود، محیط به طور خودکار همسایه‌های آن را باز نمی‌کند. این تصمیم مهم است، زیرا عامل باید هر خانه جدید را از طریق پایگاه دانش خودش به دست آورد و اطلاعات پنهان محیط به طور ناخواسته وارد حافظه عامل نشود.

\section{ساختار کلی عامل}
کد عامل در فایل
\lr{\texttt{main.py}}
قرار دارد و از دو بخش اصلی تشکیل شده است:

\begin{itemize}
    \item یک صفحه داخلی یا
          \lr{Shadow Board}
          که فقط دانسته‌های عامل را نگهداری می‌کند؛
    \item یک پایگاه دانش
          \lr{pyDatalog}
          که حقایق صفحه را دریافت می‌کند و با دو قانون اصلی، خانه‌های امن و مین‌ها را نتیجه می‌گیرد.
\end{itemize}

برای نمایش وضعیت خانه‌ها در صفحه داخلی از دو ثابت استفاده شده است:

\begin{latin}
    \begin{lstlisting}[language=Python]
AGENT_UNKNOWN = -1
AGENT_FLAGGED = -2
\end{lstlisting}
\end{latin}

عددهای صفر تا هشت نیز مقدار واقعی خانه‌های آشکارشده‌اند.

\subsection{همسایگی خانه‌ها}
تابع
\lr{\texttt{get\_neighbors}}
همه همسایه‌های معتبر یک خانه را در هشت جهت برمی‌گرداند. شرط مرزی داخل همین تابع بررسی می‌شود؛ بنابراین خانه‌های گوشه سه همسایه، خانه‌های لبه پنج همسایه و خانه‌های داخلی هشت همسایه دارند.

\begin{latin}
    \begin{lstlisting}[language=Python]
for dr in (-1, 0, 1):
    for dc in (-1, 0, 1):
        if dr == 0 and dc == 0:
            continue
        nr, nc = r + dr, c + dc
        if 0 <= nr < rows and 0 <= nc < cols:
            neighbors.append((nr, nc))
\end{lstlisting}
\end{latin}

این رابطه در شروع بازی فقط یک بار به صورت حقایق
\lr{\texttt{Adjacent(r,c,r2,c2)}}
در پایگاه دانش ثبت می‌شود و در طول بازی تغییر نمی‌کند.

\section{مدل‌سازی با منطق مرتبه اول}
محمول‌های مورد استفاده در مدل منطقی عبارت‌اند از:

\begin{itemize}
    \item \lr{\texttt{Adjacent(r,c,r2,c2)}}: خانه دوم همسایه خانه اول است.
    \item \lr{\texttt{Value(r,c,n)}}: خانه آشکارشده $(r,c)$ عدد $n$ را نشان می‌دهد.
    \item \lr{\texttt{FlaggedNeighborCount(r,c,f)}}: تعداد همسایه‌های پرچم‌گذاری‌شده برابر $f$ است.
    \item \lr{\texttt{HiddenNeighborCount(r,c,u)}}: تعداد همسایه‌های ناشناخته و بدون پرچم برابر $u$ است.
    \item \lr{\texttt{UnknownCell(r,c)}}: خانه هنوز برای عامل ناشناخته است.
    \item \lr{\texttt{Safe(r,c)}}: از حقایق فعلی نتیجه شده است که خانه امن است.
    \item \lr{\texttt{Mine(r,c)}}: از حقایق فعلی نتیجه شده است که خانه مین دارد.
\end{itemize}

در این پیاده‌سازی، حقایق همسایگی ایستا هستند، اما مقدار خانه‌ها و شمارنده‌های اطراف آن‌ها پویا هستند. این جداسازی باعث می‌شود اطلاعات هندسی صفحه فقط یک بار ساخته شود و اطلاعات وابسته به وضعیت بازی در هر دور به‌روزرسانی شوند.

\subsection{قانون امنیت}
اگر عدد یک خانه آشکارشده با تعداد همسایه‌های پرچم‌گذاری‌شده آن برابر باشد، همه همسایه‌های ناشناخته باقی‌مانده امن هستند. شکل منطقی قانون چنین است:

\[
    \begin{aligned}
         & \operatorname{Value}(r,c,n) \land
        \operatorname{FlaggedCount}(r,c,n)                        \\
         & \land\ \operatorname{Adjacent}(r,c,r',c')
        \land \operatorname{Unknown}(r',c')                       \\
         & \hspace{3cm}\Longrightarrow \operatorname{Safe}(r',c')
    \end{aligned}
\]

نسخه
\lr{pyDatalog}
همین قانون در تابع
\lr{\texttt{init\_rules}}
تعریف شده است:

\begin{latin}
    \begin{lstlisting}[language=Python]
Safe(R2, C2) <= (
    Value(R, C, N)
    & FlaggedNeighborCount(R, C, N)
    & Adjacent(R, C, R2, C2)
    & UnknownCell(R2, C2)
)
\end{lstlisting}
\end{latin}

برای مثال اگر خانه‌ای مقدار دو داشته باشد و دقیقا دو همسایه آن پرچم‌گذاری شده باشند، هیچ مین دیگری نمی‌تواند در میان همسایه‌های ناشناخته آن وجود داشته باشد.

\subsection{قانون خطر}
فرض می‌کنیم عدد خانه آشکارشده برابر $n$، تعداد پرچم‌های مجاور برابر $f$ و تعداد خانه‌های ناشناخته بدون پرچم برابر $u$ باشد. اگر

\[
    n=f+u,
\]

آنگاه همه $u$ خانه ناشناخته باید مین باشند. شکل منطقی قانون به صورت زیر است:

\[
    \begin{aligned}
         & \operatorname{Value}(r,c,n) \land
        \operatorname{FlaggedCount}(r,c,f) \land
        \operatorname{HiddenCount}(r,c,u)                         \\
         & \land\ n=f+u \land u>0
        \land \operatorname{Adjacent}(r,c,r',c')
        \land \operatorname{Unknown}(r',c')                       \\
         & \hspace{3cm}\Longrightarrow \operatorname{Mine}(r',c')
    \end{aligned}
\]

در کد، این قاعده به شکل زیر نوشته شده است:

\begin{latin}
    \begin{lstlisting}[language=Python]
Mine(R2, C2) <= (
    Value(R, C, N)
    & FlaggedNeighborCount(R, C, F)
    & HiddenNeighborCount(R, C, U)
    & (N == F + U)
    & (U > 0)
    & Adjacent(R, C, R2, C2)
    & UnknownCell(R2, C2)
)
\end{lstlisting}
\end{latin}

وجود شرط
\lr{\texttt{U > 0}}
از تولید نتیجه‌های بی‌اثر برای خانه‌هایی جلوگیری می‌کند که دیگر همسایه ناشناخته ندارند.

\section{ساخت و به‌روزرسانی پایگاه دانش}
تابع
\lr{\texttt{update\_knowledge\_base}}
در ابتدای هر دور اجرا می‌شود. ابتدا تمام حقایق پویای دور قبل با استفاده از فهرست
\lr{\texttt{dynamic\_facts}}
حذف می‌شوند. سپس صفحه داخلی پیمایش می‌شود و برای هر خانه، حقایق جدید متناسب با وضعیت فعلی ساخته می‌شوند.

برای یک خانه آشکارشده، تعداد پرچم‌ها و خانه‌های ناشناخته فقط از روی صفحه داخلی عامل محاسبه می‌شود:

\begin{latin}
    \begin{lstlisting}[language=Python]
flagged_count = sum(
    1 for nr, nc in neighbors
    if agent_board[(nr, nc)] == AGENT_FLAGGED
)
hidden_count = sum(
    1 for nr, nc in neighbors
    if agent_board[(nr, nc)] == AGENT_UNKNOWN
)
\end{lstlisting}
\end{latin}

بعد از ساخته شدن حقایق، تابع
\lr{\texttt{query\_solver}}
دو پرس‌وجوی
\lr{\texttt{Safe(R,C)}}
و
\lr{\texttt{Mine(R,C)}}
را اجرا می‌کند. پاسخ‌ها به دو فهرست مختصات تبدیل می‌شوند و نتیجه تکراری در فهرست قرار نمی‌گیرد.

\section{حلقه اصلی بازی}
مراحل عامل در هر دور به ترتیب زیر هستند:

\begin{enumerate}
    \item بازسازی حقایق پویای پایگاه دانش از روی صفحه داخلی؛
    \item پرس‌وجوی خانه‌های امن و مین‌ها؛
    \item آشکار کردن همه خانه‌هایی که امن نتیجه شده‌اند؛
    \item پرچم‌گذاری خانه‌هایی که مین نتیجه شده‌اند؛
    \item انتخاب یک حدس فقط وقتی هیچ حرکت قطعی تولید نشده باشد؛
    \item نمایش وضعیت جدید بازی و آغاز دور بعد.
\end{enumerate}

خانه‌های امن پیش از مین‌ها اعمال می‌شوند. هر عمل نیز فقط در صورتی انجام می‌شود که خانه هنوز در صفحه داخلی ناشناخته باشد. در نتیجه اگر یک پاسخ در چند دور یا از چند حقیقت مختلف دوباره نتیجه شود، عمل تکراری روی محیط اجرا نمی‌شود.

نقطه شروع در سازنده محیط آشکار شده است. از آنجا که صورت پروژه تضمین می‌کند مقدار این خانه صفر است، عامل مقدار صفر را در صفحه داخلی ثبت می‌کند و استنتاج را از همسایه‌های آن آغاز می‌کند.

\section{مدیریت بن‌بست منطقی}
دو قانون اصلی همیشه برای حل کامل همه نقشه‌ها کافی نیستند. ممکن است چند آرایش متفاوت از مین‌ها با حقایق فعلی سازگار باشند و هیچ خانه‌ای به طور قطعی امن یا مین نباشد. در این حالت تابع
\lr{\texttt{get\_safest\_guess}}
برای خانه‌های مرزی یک تخمین ساده از خطر می‌سازد.

برای هر خانه آشکارشده $i$ با عدد $n_i$، تعداد پرچم $f_i$ و تعداد همسایه ناشناخته $u_i$، خطر محلی هر همسایه ناشناخته به صورت زیر در نظر گرفته می‌شود:

\[
    p_i = \frac{n_i-f_i}{u_i}.
\]

اگر یک خانه ناشناخته $x$ کنار چند سرنخ آشکارشده باشد، میانگین تخمین‌های محلی آن‌ها محاسبه می‌شود:

\[
    \widehat{P}(\operatorname{Mine}(x))
    = \frac{1}{|C(x)|}\sum_{i\in C(x)}p_i,
\]

که $C(x)$ مجموعه سرنخ‌های مجاور خانه $x$ است. عامل خانه‌ای را انتخاب می‌کند که کمترین میانگین را داشته باشد. اگر هیچ سرنخی برای ساخت این تخمین وجود نداشته باشد، یک خانه ناشناخته به صورت تصادفی انتخاب می‌شود.

این مقدار یک احتمال دقیق مشترک نیست، زیرا وابستگی قیدهای هم‌پوشان را مدل نمی‌کند. با این حال، از حدس کاملا تصادفی بهتر از اطلاعات موجود استفاده می‌کند و فقط پس از تمام شدن حرکت‌های قطعی فعال می‌شود.

\section{آزمون قوانین منطقی}
پیش از اجرای سناریوهای کامل، دو حالت کوچک سه در سه برای بررسی مستقل قوانین ساخته شد:

\begin{enumerate}
    \item خانه مرکزی مقدار یک داشت، یکی از همسایه‌های آن پرچم‌گذاری شده بود و هفت همسایه دیگر ناشناخته بودند. پرس‌وجو دقیقا هر هفت خانه باقی‌مانده را امن برگرداند و هیچ مین جدیدی نتیجه نشد.
    \item خانه مرکزی مقدار هشت داشت و هر هشت همسایه آن ناشناخته بودند. پرس‌وجو دقیقا هر هشت همسایه را مین برگرداند و هیچ خانه امنی نتیجه نشد.
\end{enumerate}

این دو آزمون مرزی نشان می‌دهند شرط‌های شمارشی قانون امنیت و قانون خطر به همان صورتی که در مدل نوشته شده‌اند در موتور استنتاج نیز اجرا می‌شوند.

\section{نحوه ارزیابی}
چهار سناریوی مشخص‌شده در صورت پروژه اجرا شدند. برای اینکه زمان به‌دست‌آمده مربوط به استنتاج باشد، اجرای آزمایشی به صورت
\lr{headless}
انجام شد و فراخوانی نمایش گرافیکی و تاخیر ثابت
\lr{\texttt{time.sleep(0.2)}}
از اندازه‌گیری کنار گذاشته شد. خود منطق عامل، ترتیب اعمال حرکت‌ها و راهبرد حدس بدون تغییر باقی ماندند.

هر سناریو در یک فرایند جداگانه اجرا شد تا حقایق ایستای
\lr{pyDatalog}
بین اندازه‌های مختلف صفحه مشترک نشوند. در سناریوی سرعت، همان مقدار
\lr{\texttt{seed=-1}}
به محیط داده شد. فقط برای تکرارپذیر شدن همین اندازه‌گیری، مولد تصادفی پایتون پیش از ساخت صفحه با مقدار
\lr{2026}
تنظیم شد.

\section{نتایج اجرا}
\begin{center}
    \renewcommand{\arraystretch}{1.35}
    \small
    \resizebox{\textwidth}{!}{%
        \begin{tabular}{c c c c c c c c c}
            \toprule
            \textbf{ابعاد} & \textbf{مین} & \textbf{Seed} & \textbf{نتیجه}   & \textbf{پیشرفت} & \textbf{دورها} & \textbf{امن منطقی} & \textbf{مین منطقی} & \textbf{زمان} \\
            \midrule
            $9\times9$     & 9            & 99            & برد              & $100\%$         & 16             & 71                 & 9                  & $0.407\,s$    \\
            $15\times15$   & 35           & 42            & برد              & $100\%$         & 20             & 189                & 35                 & $1.042\,s$    \\
            $20\times20$   & 75           & 123           & باخت پس از 9 حدس & $94.46\%$       & 53             & 297                & 60                 & $5.537\,s$    \\
            $45\times80$   & 200          & $-1$          & برد              & $100\%$         & 43             & 3399               & 200                & $40.626\,s$   \\
            \bottomrule
        \end{tabular}%
    }
\end{center}

در سناریوی ساده و استاندارد، عامل بدون هیچ حدسی صفحه را کامل حل کرد. تعداد خانه‌های امن منطقی در این دو اجرا به ترتیب ۷۱ و ۱۸۹ است؛ با احتساب خانه شروع، این عددها دقیقا برابر تعداد کل خانه‌های امن صفحه می‌شوند. همه مین‌ها نیز به طور صحیح پرچم‌گذاری شدند.

در سناریوی چالشی، عامل ابتدا با استنتاج قطعی بخش عمده صفحه را حل کرد، اما در مجموع ۹ بار به بن‌بست رسید. حدس نهم روی یک مین قرار گرفت و بازی با پیشرفت
$94.46\%$
تمام شد. این مقدار از حداقل پیشرفت ۸۰ درصد خواسته‌شده برای این سناریو بیشتر است. همچنین همه ۶۰ پرچم قرارگرفته تا لحظه باخت صحیح بودند.

در سناریوی بزرگ، عامل همه ۳۴۰۰ خانه امن را آشکار و هر ۲۰۰ مین را پرچم‌گذاری کرد. این آزمایش نشان داد پیاده‌سازی روی صفحه
$45\times80$
بدون کرش اجرا می‌شود. زمان تقریبا ۴۰.۶ ثانیه نشان می‌دهد بازسازی حقایق پویا و پرس‌وجوی کلی موتور استنتاج در صفحه‌های بزرگ هزینه قابل توجهی دارد، هرچند اجرای کامل همچنان ممکن است.

\begin{figure}[H]
    \centering
    \includegraphics[width=0.42\textwidth]{result-9x9.png}
    \caption{وضعیت نهایی سناریوی ساده؛ پیشرفت ۱۰۰ درصد و ۹ پرچم صحیح}
\end{figure}

\begin{figure}[H]
    \centering
    \includegraphics[width=0.72\textwidth]{result-20x20.png}
    \caption{سناریوی چالشی در لحظه برخورد حدس نهم با مین؛ پیشرفت ۹۴ درصد}
\end{figure}

\section{تحلیل پیاده‌سازی}
\subsection{دلیل موفقیت قوانین اصلی}
قانون امنیت و قانون خطر دو حالت حدی یک قید شمارشی را پوشش می‌دهند. اگر همه مین‌های لازم قبلا با پرچم مشخص شده باشند، بقیه همسایه‌ها امن هستند. اگر تعداد مین باقی‌مانده دقیقا با تعداد خانه‌های ناشناخته برابر باشد، همه آن خانه‌ها مین هستند. اجرای دسته‌ای نتیجه‌ها باعث می‌شود در هر دور چندین خانه آشکار یا پرچم‌گذاری شوند و حقایق جدید برای دور بعد تولید شوند.

\subsection{حفظ مرز اطلاعات عامل و محیط}
عامل برای تصمیم‌گیری فقط از
\lr{agent\_board}
استفاده می‌کند. آرایه
\lr{\texttt{\_grid\_values}}
که جای واقعی مین‌ها را نگهداری می‌کند، در کد حل‌کننده خوانده نمی‌شود. مقدار هر خانه فقط از خروجی
\lr{\texttt{reveal}}
و وضعیت پرچم فقط از عمل خود عامل وارد حافظه می‌شود. بنابراین نتیجه‌های منطقی بر پایه اطلاعات قابل مشاهده بازی هستند.

\subsection{هزینه اجرای موتور استنتاج}
تعداد حقایق همسایگی با تعداد خانه‌ها رشد خطی دارد و برای بیشتر خانه‌های داخلی هشت رابطه ساخته می‌شود. علاوه بر آن، در هر دور کل صفحه داخلی پیمایش و همه حقایق پویا دوباره تولید می‌شوند. این طراحی ساده و قابل فهم است، اما در صفحه بزرگ، زمان هر دور بیشتر می‌شود. نتیجه سناریوی سرعت نیز همین تفاوت مقیاس را نشان می‌دهد: صفحه ۸۱ خانه‌ای در کمتر از نیم ثانیه و صفحه ۳۶۰۰ خانه‌ای در حدود ۴۰ ثانیه حل شد.

\section{محدودیت‌ها و پیشنهادهای بهبود}
مهم‌ترین محدودیت‌ها و مسیرهای بهبود نسخه فعلی عبارت‌اند از:

\begin{itemize}
    \item قوانین فعلی فقط استنتاج مستقیم پیرامون یک سرنخ را انجام می‌دهند. استنتاج مجموعه‌ای بین دو سرنخ هم‌پوشان می‌تواند بن‌بست‌های بیشتری را بدون حدس حل کند.
    \item میانگین خطر محلی یک تخمین ساده است و سازگاری هم‌زمان همه قیدها را بررسی نمی‌کند. شمارش مدل‌های سازگار یا حل مؤلفه‌های مرزی می‌تواند احتمال دقیق‌تری تولید کند.
    \item پایگاه دانش پویا در هر دور از ابتدا ساخته می‌شود. به‌روزرسانی افزایشی فقط برای خانه‌هایی که وضعیتشان تغییر کرده است می‌تواند هزینه صفحه‌های بزرگ را کاهش دهد.
    \item در صورت برابر بودن خطر چند خانه، انتخاب فعلی به ترتیب پیمایش دیکشنری وابسته است. می‌توان معیارهای تکمیلی مانند بیشترین تعداد سرنخ مجاور یا بیشترین اطلاعات مورد انتظار را برای شکستن تساوی استفاده کرد.
\end{itemize}

\end{document}
